\documentclass[11pt]{article}
\usepackage[utf8]{inputenc}
\usepackage[T1]{fontenc}
\usepackage[margin=0.85in]{geometry}
\usepackage{booktabs}
\usepackage{tabularx}
\usepackage{adjustbox}
\usepackage{graphicx}
\usepackage{amsmath}
\usepackage{amssymb}
\usepackage{hyperref}
\usepackage[expansion=false]{microtype}

\hypersetup{
    colorlinks=true,
    linkcolor=blue,
    citecolor=blue,
    urlcolor=blue
}

\title{\textbf{Proliferation Confounding and Bulk Transcriptomic Limitations of the \textit{SMARCA4}--\textit{HIF1A} Axis Across Solid Tumors}}
\author{Sachsin Sinnathurai\\
Department of Anatomy \& Cell Biology, McGill University, Montreal, QC, Canada\\
\texttt{sachsin.sinnathurai@mail.mcgill.ca}}
\date{November 2026}

\begin{document}
\maketitle

\begin{abstract}
We tested whether combined bulk mRNA expression of the chromatin remodeler \textit{SMARCA4} and the hypoxia factor \textit{HIF1A} provides proliferation-independent prognostic or pharmacogenomic stratification across adult solid tumors, and whether normoxic cell culture alters \textit{HIF1A}--\textit{SLC7A1} co-expression relative to primary patient tumors. Across an 18-cohort adult solid TCGA panel (16 stage- or grade-adjusted; \texttt{SARC} and \texttt{GBM} age-adjusted; $N = 7,382$, $2,352$ deaths), an unweighted composite of standardized \textit{SMARCA4} and \textit{HIF1A} mRNA did not predict overall survival in DerSimonian-Laird random-effects meta-analysis ($\text{HR}^* = 1.044$, $95\%$ CI $[0.983\text{--}1.108]$, $p = 0.162$, $I^2 = 49.4\%$) and attenuated further after \textit{MKI67} proliferation adjustment ($\text{HR}^* = 1.018$, $95\%$ CI $[0.962\text{--}1.077]$, $p = 0.545$, $I^2 = 37.2\%$). Adding un-staged \texttt{TCGA-PRAD} ($N = 493$, $10$ deaths) in a 19-cohort sensitivity pool ($N = 7,875$) yielded the same null result (base $\text{HR}^* = 1.048$, $p = 0.130$; +\textit{MKI67} $\text{HR}^* = 1.020$, $p = 0.488$). Under Benjamini-Hochberg false discovery rate (BH-FDR) control after \textit{MKI67} adjustment, $0/18$ primary cohorts ($0/19$ in the sensitivity family) remained significant at $q < 0.10$ (\texttt{TCGA-KIRC} $q_{18} = 0.125$, \texttt{TCGA-BLCA} $q_{18} = 0.164$). By contrast, a 14-gene non-cell-cycle hypoxia metagene (\textit{Buffa-14}, excluding \textit{CDKN3}) remained prognostic across the 18-cohort TCGA panel ($\text{HR}^* = 1.132$, $p = 4.26 \times 10^{-6}$; +\textit{MKI67} $\text{HR}^* = 1.115$, $p = 7.80 \times 10^{-7}$, $I^2 = 0.9\%$) and improved omnibus nested model fit ($\chi^2(18) = 56.10$, $p = 8.60 \times 10^{-6}$; +\textit{MKI67} $\chi^2(18) = 42.67$, $p = 8.90 \times 10^{-4}$). Across Sanger GDSC2 screens ($m = 286$ drugs) on empirical $1 - \text{AUC}$, $0/286$ compounds cleared $10\%$ BH-FDR after \textit{MKI67} adjustment ($\min q = 0.1126$). Finally, when Broad DepMap 24Q2 cell lines were restricted to the three lineages matching our primary TCGA expression panel (Breast, CNS/Brain, Pancreas), \textit{HIF1A}--\textit{SLC7A1} co-expression did not differ significantly from primary tumors ($n = 229$ unique models: $r = +0.2154$, Fisher $Z = 1.02$, $p = 0.305$ vs.\ TCGA $r = +0.2844$). Bulk \textit{SMARCA4} and \textit{HIF1A} mRNA thus do not form a generalizable solid-tumor prognostic or therapeutic biomarker once proliferation and lineage composition are controlled.
\end{abstract}

\section{Introduction}
While hypoxia-inducible transcription factors (HIF-1$\alpha$ and HIF-2$\alpha$) are regulated post-translationally via oxygen-dependent prolyl hydroxylases (PHD1--3) and von Hippel-Lindau (VHL) ubiquitin-mediated degradation~\cite{ivan2001,jaakkola2001}, chromatin-remodeling cooperation by the mammalian SWI/SNF (BAF) ATPase subunit \textit{SMARCA4} (BRG1)~\cite{kenneth2009,sena2013} and hypoxia-associated metabolic and drug-efflux programs raise the question of whether bulk \textit{SMARCA4} and \textit{HIF1A} mRNA abundance carries proliferation-independent prognostic or pharmacogenomic utility across solid tumors. Specifically, HIF-2$\alpha$ (\textit{EPAS1}) drives expression of the neutral amino acid carrier \textit{SLC7A5}~\cite{elorza2012} and a mitochondrial glutamine-transporter isoform of \textit{SLC1A5} (\texttt{SLC1A5\_var})~\cite{yoo2020}, while standard bulk RNA-seq gene counts aggregate canonical plasma-membrane \textit{SLC1A5} with mitochondrial \texttt{SLC1A5\_var}. Alongside \textit{SLC1A5}, we examined the cationic amino acid transporter \textit{SLC7A1}---which shares CRISPR co-dependency with \textit{SLC1A5} across solid tumor cell lines in Broad DepMap ($r = +0.207$)---and evaluated the multidrug efflux pump \textit{ABCB1} as an exploratory drug-efflux comparator in PDAC single-cell data.

Testing whether bulk \textit{SMARCA4} and \textit{HIF1A} mRNA levels provide generalizable clinical or pharmacogenomic stratification requires explicit control for three confounders. First, because oxygen tension regulates HIF-1$\alpha$ post-translationally, steady-state \textit{HIF1A} transcript abundance can diverge from downstream hypoxic pathway output. Second, rapidly proliferating tumors co-upregulate chromatin-remodeling genes alongside \textit{MKI67}, creating proliferation-driven survival and drug-response associations unless proliferation is modeled directly. Third, expression comparisons across mismatched tumor lineages between patient cohorts and cell-line panels can introduce apparent differences unrelated to ambient $21\%\ \text{O}_2$ oxygen tension.

Here we evaluate the unweighted \textit{SMARCA4} + \textit{HIF1A} mRNA composite alongside \textit{MKI67} proliferation and the \textit{Buffa-14} hypoxia signature~\cite{buffa2010} across the 18-cohort TCGA panel ($N = 7,382$, with 19-cohort \texttt{TCGA-PRAD} sensitivity), three independent non-TCGA clinical cohorts, Broad DepMap 24Q2 functional genomics, four $1\%$ vs.\ $21\%\ \text{O}_2$ hypoxia series~\cite{gse139673,gse145108,gse283995,gse300765}, Sanger GDSC2 drug screens, pancreatic ductal adenocarcinoma (PDAC) single-cell RNA-seq, and TCGA proteomic and genomic scar datasets.

\section{Results}

\subsection{Survival Across the 18-Cohort TCGA Panel Does Not Stratify Beyond Proliferation and \textit{Buffa-14}}
We evaluated the unweighted composite score $Z(z_{\textit{SMARCA4}} + z_{\textit{HIF1A}})$ across the 18-cohort TCGA panel ($N = 7,382$, $2,352$ deaths)~\cite{tcga2018}, treating the 19-cohort pool including un-staged \texttt{TCGA-PRAD} ($N = 493$, $10$ deaths) as a sensitivity check ($N = 7,875$, $2,362$ deaths). Downstream metabolic and efflux genes (\textit{SLC7A1}, \textit{SLC1A5}, and \textit{ABCB1}) were examined separately in DepMap functional dependency, hypoxic perturbation, and PDAC single-cell datasets. Across the 18-cohort TCGA panel using DerSimonian-Laird random-effects meta-analysis, the \textit{SMARCA4} + \textit{HIF1A} score did not predict overall survival ($\text{HR}^* = 1.044$, $95\%$ CI $[0.983\text{--}1.108]$, $p = 0.162$, $I^2 = 49.4\%$; Table~\ref{tab:survival}, Figure~\ref{fig:summary}A). Including \textit{MKI67} to adjust for tumor proliferation reduced the pooled hazard to $\text{HR}^* = 1.018$ ($95\%$ CI $[0.962\text{--}1.077]$, $p = 0.545$, $I^2 = 37.2\%$). To verify that purity adjustment in \texttt{TCGA-PAAD} did not distort pooling, replacing purity-adjusted \texttt{TCGA-PAAD} ($N = 168$) with a uniform \texttt{Age + Stage/Grade} model (\texttt{TCGA-PAAD} $N = 174$: base $\text{HR} = 1.081$, $p = 0.501$; +\textit{MKI67} $\text{HR} = 1.000$, $p = 0.997$) yielded virtually identical 18-cohort meta-estimates ($N = 7,388$: base $\text{HR}^* = 1.041$, $95\%$ CI $[0.981\text{--}1.105]$, $p = 0.180$; +\textit{MKI67} $\text{HR}^* = 1.016$, $95\%$ CI $[0.961\text{--}1.075]$, $p = 0.580$). The 19-cohort sensitivity pool gave matching null results at baseline ($\text{HR}^* = 1.048$, $95\%$ CI $[0.986\text{--}1.114]$, $p = 0.130$, $I^2 = 50.5\%$) and after \textit{MKI67} control ($\text{HR}^* = 1.020$, $95\%$ CI $[0.964\text{--}1.080]$, $p = 0.488$, $I^2 = 36.8\%$). Across the 11 remaining cohorts in the 19-cohort sensitivity pool ($N = 4,552$, i.e., $7,875$ minus the 8 TCGA cohorts selected for secondary sensitivity modeling [$N = 3,323$], including un-staged \texttt{TCGA-PRAD}), the pooled hazard centered near $\text{HR} = 1.0$ ($\text{HR}^* = 0.962$, $p = 0.186$; 10-cohort remainder excluding \texttt{TCGA-PRAD}: $N = 4,059$, $\text{HR}^* = 0.958$, $p = 0.139$).

Under BH-FDR control across the 18-cohort TCGA panel (Table~\ref{tab:survival}), only clear cell renal carcinoma (\texttt{TCGA-KIRC}, $N = 531$, $\text{HR} = 1.289$, $p = 6.60 \times 10^{-4}$, $q_{18} = 0.012$) cleared $5\%$ or $10\%$ BH-FDR at baseline without \textit{MKI67}, whereas urothelial bladder (\texttt{TCGA-BLCA}, $q_{18} = 0.211$) and breast (\texttt{TCGA-BRCA}, $q_{18} = 0.211$) did not. After \textit{MKI67} adjustment, $0/18$ cohorts cleared $5\%$ or $10\%$ BH-FDR (\texttt{TCGA-KIRC} $\text{HR} = 1.239$, $p = 0.007$, $q_{18} = 0.125$; \texttt{TCGA-BLCA} $\text{HR} = 1.230$, $p = 0.018$, $q_{18} = 0.164$; \texttt{TCGA-UCEC} $\text{HR} = 0.831$, $p = 0.059$, $q_{18} = 0.355$). Enforcing BH-FDR across all 19 cohorts sensitivity-wide gave the same null result (\texttt{TCGA-KIRC} $q_{19} = 0.132$; \texttt{TCGA-BLCA} $q_{19} = 0.173$; $0/19$ at $q < 0.10$).

To test whether poor surrogacy of bulk \textit{HIF1A} mRNA explained this null result, we benchmarked the 14-gene non-cell-cycle \textit{Buffa-14} hypoxia signature (excluding \textit{CDKN3}~\cite{buffa2010}) across the 18-cohort TCGA panel ($N = 7,382$). Unlike the two-gene score, \textit{Buffa-14} alone stratified mortality at baseline ($\text{HR}^* = 1.132$, $95\%$ CI $[1.074\text{--}1.194]$, $p = 4.26 \times 10^{-6}$, $I^2 = 34.2\%$) and persisted after \textit{MKI67} adjustment ($\text{HR}^* = 1.115$, $95\%$ CI $[1.068\text{--}1.165]$, $p = 7.80 \times 10^{-7}$, $I^2 = 0.9\%$). In omnibus summed nested likelihood-ratio tests across the 18 cohorts ($\chi^2(18)$), adding \textit{Buffa-14} onto \textit{SMARCA4} + \textit{HIF1A} improved model fit at baseline ($\chi^2(18) = 56.10$, $p = 8.60 \times 10^{-6}$) and after \textit{MKI67} control ($\chi^2(18) = 42.67$, $p = 8.90 \times 10^{-4}$), testing whether \textit{Buffa-14} improves model fit in at least one cohort rather than uniformly across every cancer type. Results were identical across the 19-cohort sensitivity pool (\textit{Buffa-14} base $\text{HR}^* = 1.133$, $p = 2.17 \times 10^{-6}$, $\chi^2(19) = 56.56$, $p = 1.33 \times 10^{-5}$; +\textit{MKI67} $\text{HR}^* = 1.116$, $p = 5.54 \times 10^{-7}$, $\chi^2(19) = 43.03$, $p = 1.29 \times 10^{-3}$), whereas adding \textit{SMARCA4} + \textit{HIF1A} onto \textit{Buffa-14} + \textit{MKI67} failed to improve fit in either the 18-cohort panel ($\chi^2(18) = 27.90$, $p = 0.064$) or the 19-cohort sensitivity pool ($\chi^2(19) = 29.34$, $p = 0.061$).

We next evaluated three independent non-TCGA clinical cohorts (\texttt{CGGA-325}~\cite{cgga2021}, \texttt{BRCA-METABRIC}~\cite{curtis2012}, and \texttt{ICGC/PCAWG-PAAD}~\cite{pcawg2020}; Table~\ref{tab:survival}) alongside subtype and purity sensitivity checks within individual TCGA cohorts. In diffuse glioma (\texttt{CGGA-325}, $N = 313$, $218$ deaths), adding \textit{MKI67} completely abolished survival stratification (baseline $\text{HR} = 1.193$, $p = 0.007 \rightarrow +\textit{MKI67}$ $\text{HR} = 1.006$, $p = 0.943$; \textit{TOP2A}: $\text{HR} = 0.949$, $p = 0.505$). In breast cancer, adjusting for PAM50 subtype rendered \texttt{TCGA-BRCA} null ($N = 950$, $\text{HR} = 1.052$, $p = 0.611$), and \texttt{BRCA-METABRIC} ($N = 1,466$, $823$ deaths) attenuated progressively from baseline ($\text{HR} = 1.114$, $p = 0.002$) through +\textit{MKI67} ($\text{HR} = 1.085$, $p = 0.025$) to full subtype and receptor adjustment (\texttt{+PAM50/Claudin-low + ER/PR/HER2}, $N = 1,463$, $\text{HR} = 1.035$, $p = 0.361$). In \texttt{TCGA-KIRC} ($N = 496$) and \texttt{TCGA-BLCA} ($N = 393$) with ABSOLUTE purity data~\cite{carter2012}, nominal associations persisted after +\textit{MKI67}+Purity adjustment (\texttt{TCGA-KIRC} $\text{HR} = 1.282$, $p = 0.002$; \texttt{TCGA-BLCA} $\text{HR} = 1.214$, $p = 0.028$), though neither cleared 18-cohort BH-FDR. Single-gene decomposition across the 10-cohort clinical subset ($N = 5,102$: \texttt{TCGA-KIRC}, \texttt{TCGA-BLCA}, \texttt{TCGA-GBM}, \texttt{TCGA-BRCA}, \texttt{TCGA-SARC}, \texttt{TCGA-SKCM}, \texttt{TCGA-LIHC}, and \texttt{TCGA-PAAD} [$N = 3,323$], plus \texttt{CGGA-325} and \texttt{BRCA-METABRIC}) showed that any residual unadjusted signal derived from \textit{SMARCA4} ($\text{HR}^* = 1.153$, $p = 2.82 \times 10^{-4}$), whereas \textit{HIF1A} mRNA alone was null ($\text{HR}^* = 1.047$, $p = 0.222$).

In pancreatic cancer, bulk \texttt{TCGA-PAAD} survival was null across complete-case ($N = 168$, $\text{HR} = 1.143$, $p = 0.273$, $C$-index $= 0.579$; +\textit{MKI67} $\text{HR} = 1.039$, $p = 0.749$), Raphael et al.\ histologically restricted ductal PDAC ($N = 148$, $\text{HR} = 0.950$, $p = 0.669$, $C$-index $= 0.578$~\cite{raphael2017}), and cBioPortal pure-ductal \texttt{Age+Stage+Purity} sensitivity modeling ($N = 138$, $\text{HR} = 0.996$, $p = 0.974$; $p = 0.273\text{--}0.974$ across sub-models). Although the smaller independent \texttt{ICGC/PCAWG-PAAD} cohort~\cite{pcawg2020} ($N = 71$, $48$ deaths; stage and purity unavailable) reached significance after age and sex adjustment ($\text{HR} = 1.504$, $95\%$ CI $[1.109\text{--}2.041]$, $p = 0.009$, $C$-index $= 0.658$), unadjusted log-rank separation was borderline ($p = 0.060$). The discrepancy between \texttt{TCGA-PAAD} (null) and \texttt{ICGC/PCAWG-PAAD} ($p = 0.009$, unadjusted for stage or purity) cannot be resolved from bulk RNA-seq alone.

\begin{table*}[htbp]
\centering
\caption{\textbf{Complete 18-Cohort TCGA Survival Panel, Meta-Analyses, and Independent Non-TCGA Cohorts.} \textsuperscript{\dag}Cohort included in the 10-cohort single-gene decomposition subset ($N = 5,102$: eight TCGA cohorts [$N = 3,323$] plus \texttt{CGGA-325} and \texttt{BRCA-METABRIC}).}
\label{tab:survival}
\scriptsize
\setlength{\tabcolsep}{3pt}
\begin{adjustbox}{max width=\textwidth}
\begin{tabular}{@{}lcclccc@{}}
\toprule
\textbf{Cohort / Meta-Analysis} & \textbf{$N$} & \textbf{Deaths} & \textbf{Covariates Adjusted} & \textbf{Base HR [95\% CI]} & \textbf{Base $p$ ($q_{18}$)} & \textbf{+\textit{MKI67} HR ($p$; $q_{18}$)} \\
\midrule
TCGA-KIRC\textsuperscript{\dag} & 531 & 175 & Age, Stage & 1.289 [1.114--1.492] & $6.60\times 10^{-4}$ (0.012) & 1.239 (0.007; 0.125) \\
TCGA-BLCA\textsuperscript{\dag} & 403 & 177 & Age, Stage & 1.206 [1.024--1.419] & 0.024 (0.211) & 1.230 (0.018; 0.164) \\
TCGA-UCEC & 523 & 87 & Age, Grade & 0.832 [0.686--1.008] & 0.061 (0.274) & 0.831 (0.059; 0.355) \\
TCGA-GBM\textsuperscript{\dag} & 158 & 127 & Age & 0.922 [0.778--1.093] & 0.348 (0.528) & 0.879 (0.171; 0.565) \\
TCGA-LUAD & 485 & 177 & Age, Stage & 0.975 [0.832--1.141] & 0.749 (0.899) & 0.900 (0.186; 0.565) \\
TCGA-BRCA\textsuperscript{\dag} & 1,050 & 143 & Age, Stage & 1.187 [1.012--1.393] & 0.035 (0.211) & 1.125 (0.188; 0.565) \\
TCGA-SARC\textsuperscript{\dag} & 253 & 96 & Age & 1.183 [0.963--1.453] & 0.110 (0.329) & 1.122 (0.306; 0.630) \\
TCGA-CESC & 254 & 58 & Age, Grade & 1.142 [0.874--1.491] & 0.331 (0.528) & 1.150 (0.310; 0.630) \\
TCGA-LUSC & 469 & 200 & Age, Stage & 0.914 [0.790--1.057] & 0.226 (0.508) & 0.924 (0.339; 0.630) \\
TCGA-ESCA & 159 & 64 & Age, Stage & 0.884 [0.683--1.145] & 0.352 (0.528) & 0.883 (0.350; 0.630) \\
TCGA-HNSC & 444 & 188 & Age, Stage & 0.959 [0.831--1.107] & 0.571 (0.790) & 0.943 (0.433; 0.657) \\
TCGA-SKCM\textsuperscript{\dag} & 417 & 196 & Age, Stage & 1.106 [0.954--1.281] & 0.181 (0.465) & 1.063 (0.438; 0.657) \\
TCGA-STAD & 377 & 149 & Age, Stage & 0.971 [0.828--1.138] & 0.716 (0.899) & 1.052 (0.560; 0.776) \\
TCGA-LIHC\textsuperscript{\dag} & 343 & 118 & Age, Stage & 1.172 [0.977--1.407] & 0.088 (0.316) & 1.052 (0.607; 0.780) \\
TCGA-PAAD\textsuperscript{\dag} & 168 & 89 & Age, Stage, Purity & 1.143 [0.900--1.452] & 0.273 (0.528) & 1.039 (0.749; 0.871) \\
TCGA-THCA & 494 & 16 & Age, Stage & 0.992 [0.576--1.710] & 0.977 (0.977) & 1.088 (0.774; 0.871) \\
TCGA-COADREAD & 554 & 110 & Age, Stage & 0.986 [0.814--1.193] & 0.882 (0.934) & 0.980 (0.858; 0.908) \\
TCGA-OV & 300 & 182 & Age, Stage & 1.013 [0.873--1.175] & 0.864 (0.934) & 0.992 (0.916; 0.916) \\
\midrule
\textbf{18-Cohort TCGA Meta (\textit{SMARCA4}+\textit{HIF1A})} & 7,382 & 2,352 & Base / +\textit{MKI67} & 1.044 [0.983--1.108] & 0.162 ($I^2=49.4\%$) & 1.018 [0.962--1.077] ($p=0.545$, $I^2=37.2\%$) \\
\textbf{18-Cohort TCGA Meta (\textit{Buffa-14})} & 7,382 & 2,352 & Base / +\textit{MKI67} & 1.132 [1.074--1.194] & $4.26\times 10^{-6}$ ($I^2=34.2\%$) & 1.115 [1.068--1.165] ($p=7.80\times 10^{-7}$, $I^2=0.9\%$) \\
\textbf{19-Cohort Sensitivity (\textit{SMARCA4}+\textit{HIF1A})} & 7,875 & 2,362 & Incl.\ un-staged PRAD & 1.048 [0.986--1.114] & 0.130 ($I^2=50.5\%$) & 1.020 [0.964--1.080] ($p=0.488$, $I^2=36.8\%$) \\
\textbf{19-Cohort Sensitivity (\textit{Buffa-14})} & 7,875 & 2,362 & Incl.\ un-staged PRAD & 1.133 [1.076--1.193] & $2.17\times 10^{-6}$ ($I^2=30.7\%$) & 1.116 [1.069--1.165] ($p=5.54\times 10^{-7}$, $I^2=0.0\%$) \\
\midrule
\multicolumn{7}{@{}l}{\textit{Independent Non-TCGA Cohorts}} \\
CGGA-325 & 313 & 218 & Age, Grade, IDH, MGMT & 1.193 [1.050--1.355] & 0.007 & 1.006 ($p=0.943$; +\textit{TOP2A} $p=0.505$) \\
BRCA-METABRIC & 1,466 & 823 & Age, Stage & 1.114 [1.040--1.192] & 0.002 & 1.085 ($p=0.025$); +Receptors $p=0.361$ \\
ICGC/PCAWG-PAAD & 71 & 48 & Age, Sex & 1.504 [1.109--2.041] & 0.009 & Stage unavail.; unadj.\ log-rank $p=0.060$ \\
\bottomrule
\end{tabular}
\end{adjustbox}
\end{table*}

\subsection{Lineage-Matched DepMap Analysis Shows Preserved \textit{HIF1A}--\textit{SLC7A1} Coupling and Clarifies \textit{VHL}-LoF \textit{EPAS1} Biology}
We evaluated \textit{HIF1A}--\textit{SLC7A1} co-expression across \texttt{TCGA-BRCA} ($n = 1,050$), \texttt{TCGA-PAAD} ($n = 176$), and \texttt{TCGA-GBM} ($n = 158$; $N = 1,384$ total) and matched them to their corresponding DepMap lineages, reporting per-lineage correlations alongside the pooled estimate. Comparing these three primary TCGA cohorts (within-cohort Z-scored \textit{HIF1A}--\textit{SLC7A1} $r = +0.2844$, $p = 3.72 \times 10^{-27}$) against all 26 solid lineages in Broad DepMap 24Q2 ($n = 1,349$, within-lineage Z-scored $r = +0.0249$, $p = 0.360$) introduces a major lineage mismatch (Fisher $Z = 6.98$, $p = 2.87 \times 10^{-12}$). The \texttt{TCGA-PAAD} expression panel includes 8 tumors beyond the $N = 168$ multivariable survival set that lacked complete stage or ABSOLUTE purity covariates. Within those three TCGA cohorts, \textit{HIF1A}--\textit{SLC7A1} correlation was positive in \texttt{TCGA-BRCA} ($r = +0.3269$, $p = 1.43 \times 10^{-27}$) and \texttt{TCGA-GBM} ($r = +0.2066$, $p = 0.009$), yet null in \texttt{TCGA-PAAD} ($r = +0.0991$, $p = 0.190$). When Broad DepMap 24Q2 cell lines were restricted to the three lineages matching the TCGA expression panel (Breast $n = 71$, $r = +0.3899$; CNS/Brain $n = 103$, $r = +0.0692$; Pancreas $n = 55$, $r = +0.2655$), within-lineage Z-scored \textit{HIF1A}--\textit{SLC7A1} co-expression did not differ significantly from primary tumors across all unique matched models ($n = 229$, $r = +0.2154$, $p = 0.0010$, raw $r = +0.2504$; Fisher $Z = 1.02$, $p = 0.305$ vs.\ TCGA) or across solid-filtered models ($n = 221$, $r = +0.1736$, $p = 0.0097$, Fisher $Z = 1.61$, $p = 0.108$). Direct breast-only comparison showed slightly higher correlation in DepMap Breast lines than in \texttt{TCGA-BRCA} tumors (Fisher $Z = -0.58$, $p = 0.563$). Ambient $21\%\ \text{O}_2$ culture therefore does not ablate \textit{HIF1A}--\textit{SLC7A1} mRNA co-expression once lineages are matched.

Within DepMap 24Q2 CRISPR-Cas9 Chronos profiles~\cite{depmap2021} ($N = 1,038$ solid tumor lines), in \textit{VHL}-LoF lines ($n = 19$ vs.\ $1,019$ WT solid lines), \textit{EPAS1} exhibits selective negative Chronos dependency (mean $-0.3464 \pm 0.3966$ vs.\ $+0.0045 \pm 0.1230$ in WT, $d = -2.644$, $p = 0.001$; vs.\ $15$ WT kidney lines $+0.0451$, $p = 6.74 \times 10^{-4}$), while \textit{HIF1A} exhibits positive/non-essential Chronos scores (mean $+0.3690 \pm 0.2683$ vs.\ $+0.1127 \pm 0.1479$ in WT, Welch $p = 5.86 \times 10^{-4}$). Moreover, \textit{EPAS1} Chronos was uncorrelated with \textit{SLC7A1} ($r = -0.0877$, $p = 0.721$). The correlation between \textit{HIF1A} and \textit{SLC7A1} Chronos scores ($r = +0.5412$, $p = 0.017$) therefore reflects co-variation among non-essential fitness scores rather than lethal hypoxic dependency. Across all $N = 1,038$ solid tumor lines without extreme-group truncation, \textit{SLC1A5} and \textit{SLC7A1} Chronos scores correlated modestly ($r = +0.2074$, $\rho = +0.1485$, $p = 1.51 \times 10^{-11}$), and continuous \textit{SMARCA4} co-essentiality across all $18,434$ genes recovered core BAF subunits (\textit{SMARCB1} $r = +0.2952$, $q = 2.40 \times 10^{-18}$; \textit{SMARCC1} $r = +0.3018$, $q = 4.84 \times 10^{-19}$; \textit{SMARCD1} $r = +0.2816$, $q = 1.37 \times 10^{-16}$; \textit{ARID1A} $r = +0.2257$, $q = 8.71 \times 10^{-10}$; \textit{SMARCE1} $r = +0.2100$, $q = 2.19 \times 10^{-8}$) and the transcriptional elongation factor \textit{MLLT1} ($r = +0.1711$, $q = 2.24 \times 10^{-5}$) alongside \textit{SMARCA2} paralog buffering ($r = -0.1967$, $q = 2.53 \times 10^{-7}$). Comparing non-malignant diploid lines ($n = 16$, mean \textit{SMARCA4} $-0.0866$) against solid tumor lines ($n = 1,038$, mean $-0.3466$) gave Mann-Whitney $p = 0.007$; adjusting for \textit{MKI67} mRNA expression in OLS ($n = 14$ diploid and $974$ solid lines with RNA-seq) kept a positive difference ($\beta = +0.2502$, $p = 0.0060$).

Direct $1\%$ vs.\ $21\%\ \text{O}_2$ hypoxia experiments across four NCBI GEO series~\cite{gse139673,gse145108,gse283995,gse300765} further explain why bulk \textit{HIF1A} mRNA fails as a hypoxia surrogate (exploratory Table~\ref{tab:geo}, with BH-FDR across the 24 powered single-lineage $t$-tests; \texttt{GSE139673}~\cite{gse139673} was analyzed at gene-level abundance rather than alternative splice-isoform resolution). While \textit{CA9} ($\log_2\text{FC} = +3.079$ to $+4.095$, all raw $p \le 9.54 \times 10^{-4}$, BH $q \le 0.003$ in powered single-lineage series; $+3.977$ to $+5.776$ descriptively in \texttt{GSE145108}) and \textit{VEGFA} ($+0.775$ to $+1.813$, all raw $p \le 0.009$, BH $q \le 0.020$; $+1.116$ to $+1.423$ in \texttt{GSE145108}) rose consistently in every dataset, \textit{HIF1A} mRNA never increased ($\log_2\text{FC} = -1.589$ to $+0.107$). Neither \textit{SLC7A1} nor \textit{SLC1A5} induced consistently under hypoxia: \textit{SLC7A1} declined in AsPC-1 ($-0.458$, $p = 0.001, q = 0.003$), stayed flat in U87MG ($p = 0.517, q = 0.620$), and showed only a transient 24~h rise in HepG2 ($+0.299$, $p = 4.78 \times 10^{-5}, q = 2.30 \times 10^{-4}$; \textit{SLC1A5} $+0.128$, $p = 0.008, q = 0.020$) that vanished by day 6 (\textit{SLC7A1} $-0.200$, $p = 0.121, q = 0.182$; \textit{SLC1A5} $+0.107$, $p = 0.382, q = 0.483$). Because \texttt{GSE145108} contains only $n = 2$ pairs per distinct cancer lineage (RKO colorectal, A549 lung, and H460 lung), we report it strictly descriptively without cross-organ pooling; all three cell lines showed negative \textit{SLC7A1} ($\log_2\text{FC} = -0.254$, $-0.716$, $-1.367$) and \textit{SLC1A5} ($-0.540$, $-0.946$, $-1.298$) shifts under hypoxia.

\begin{table*}[htbp]
\centering
\caption{\textbf{Exploratory Paired $1\%$ vs.\ $21\%\ \text{O}_2$ Hypoxia Perturbation Across Four NCBI GEO Solid-Tumor Series (BH-FDR Across 24 Powered Single-Lineage Tests).} \texttt{GSE139673}~\cite{gse139673} was analyzed at gene-level abundance rather than alternative splice-isoform resolution. \texttt{---} denotes descriptive $\log_2\text{FC}$ only for \texttt{GSE145108} ($n = 2$ pairs per lineage; formal $t$-tests omitted as underpowered and cross-organ pooling avoided).}
\label{tab:geo}
\scriptsize
\setlength{\tabcolsep}{3.5pt}
\begin{adjustbox}{max width=\textwidth}
\begin{tabular}{@{}lcccccc@{}}
\toprule
\textbf{Series (Cell Model / Duration)} & \textbf{\textit{CA9} $\log_2\text{FC}$ ($p$; $q$)} & \textbf{\textit{VEGFA} $\log_2\text{FC}$ ($p$; $q$)} & \textbf{\textit{HIF1A} $\log_2\text{FC}$ ($p$; $q$)} & \textbf{\textit{SLC7A1} $\log_2\text{FC}$ ($p$; $q$)} & \textbf{\textit{SLC1A5} $\log_2\text{FC}$ ($p$; $q$)} & \textbf{\textit{SMARCA4} $\log_2\text{FC}$ ($p$; $q$)} \\
\midrule
\texttt{GSE139673}~\cite{gse139673} (AsPC-1, 24~h, $n=3$ pairs) & $+3.079$ ($5.08 \times 10^{-4}$; $0.002$) & $+1.079$ ($0.009$; $0.020$) & $+0.013$ ($0.698$; $0.736$) & $-0.458$ ($0.001$; $0.003$) & $+0.038$ ($0.736$; $0.736$) & $-0.235$ ($0.053$; $0.085$) \\
\texttt{GSE145108}~\cite{gse145108} (RKO CRC, 24~h, $n=2$ pairs) & $+3.977$ (---) & $+1.215$ (---) & $-1.589$ (---) & $-0.254$ (---) & $-0.540$ (---) & $-0.367$ (---) \\
\texttt{GSE145108}~\cite{gse145108} (A549 NSCLC, 24~h, $n=2$ pairs) & $+5.192$ (---) & $+1.423$ (---) & $-1.252$ (---) & $-0.716$ (---) & $-0.946$ (---) & $-0.434$ (---) \\
\texttt{GSE145108}~\cite{gse145108} (H460 NSCLC, 24~h, $n=2$ pairs) & $+5.776$ (---) & $+1.116$ (---) & $-0.491$ (---) & $-1.367$ (---) & $-1.298$ (---) & $-0.889$ (---) \\
\texttt{GSE283995}~\cite{gse283995} (HepG2, 24~h, $n=6$/arm) & $+3.267$ ($4.11 \times 10^{-11}$; $9.86 \times 10^{-10}$) & $+0.968$ ($8.54 \times 10^{-9}$; $1.02 \times 10^{-7}$) & $-0.083$ ($0.266$; $0.355$) & $+0.299$ ($4.78 \times 10^{-5}$; $2.30 \times 10^{-4}$) & $+0.128$ ($0.008$; $0.020$) & $+0.067$ ($0.214$; $0.302$) \\
\texttt{GSE283995}~\cite{gse283995} (HepG2, day 6, $n=5$ vs.\ 6) & $+3.376$ ($4.99 \times 10^{-6}$; $3.99 \times 10^{-5}$) & $+0.775$ ($3.75 \times 10^{-5}$; $2.25 \times 10^{-4}$) & $-0.229$ ($0.018$; $0.032$) & $-0.200$ ($0.121$; $0.182$) & $+0.107$ ($0.382$; $0.483$) & $+0.356$ ($0.011$; $0.022$) \\
\texttt{GSE300765}~\cite{gse300765} (U87MG, 48~h, $n=3$ pairs) & $+4.095$ ($9.54 \times 10^{-4}$; $0.003$) & $+1.813$ ($6.95 \times 10^{-4}$; $0.002$) & $+0.107$ ($0.573$; $0.655$) & $-0.053$ ($0.517$; $0.620$) & $-0.046$ ($0.734$; $0.736$) & $-0.376$ ($0.015$; $0.028$) \\
\bottomrule
\end{tabular}
\end{adjustbox}
\end{table*}

\subsection{Normoxic GDSC2 Pharmacology Fails Multiple-Testing Control Post-\textit{MKI67}, and Multi-Omics Associations Are Weak}
Across $m = 286$ compounds in Sanger GDSC2~\cite{iorio2016} evaluated on bounded empirical $1 - \text{AUC}$ with lineage Type II ANOVA, baseline effect sizes were uniformly small (median partial $R^2$, $\eta_p^2 = 0.176\%$; baseline screen-wide maximum $3.53\%$). Before proliferation adjustment, only Navitoclax cleared $5\%$ BH-FDR ($F(1,570) = 14.32$, $\eta_p^2 = 2.45\%$, $p = 1.71 \times 10^{-4}$, $q = 0.049$) and nine compounds cleared $10\%$ BH-FDR, while secondary $\ln(\text{IC}_{50})$ testing yielded $0/286$ at $q < 0.10$ ($\min q = 0.110$). After \textit{MKI67} adjustment, $0/286$ compounds cleared $10\%$ BH-FDR (Figure~\ref{fig:summary}B); the five top-ranked compounds (\texttt{WEHI-539}, \texttt{Navitoclax}, \texttt{Cisplatin}, \texttt{Elephantin}, and \texttt{ABT737}; raw $p = 6.43 \times 10^{-4}$ to $1.97 \times 10^{-3}$) tied at a step-up plateau of $q = 0.1126$, while the highest post-\textit{MKI67} partial $R^2$ occurred for \texttt{SL0101} ($\eta_p^2 = 4.76\%$, $p = 0.038$, $q = 0.340$). Bulk \textit{SMARCA4} + \textit{HIF1A} expression therefore does not predict selective drug response once lineage and proliferation are controlled.

\begin{figure*}[htbp]
\centering
\includegraphics[width=\textwidth]{../figures/fig1_audit_summary.pdf}
\caption{\textbf{Primary 18-cohort TCGA survival meta-analysis and Sanger GDSC2 post-\textit{MKI67} pharmacogenomic screen.} \textbf{(A)} DerSimonian-Laird random-effects per-SD hazard ratios and $95\%$ confidence intervals across the 18-cohort TCGA panel ($N = 7,382$, $2,352$ deaths) comparing the unweighted \textit{SMARCA4} + \textit{HIF1A} score (base $\text{HR}^* = 1.044, p = 0.162$; +\textit{MKI67} $\text{HR}^* = 1.018, p = 0.545$) against the non-cell-cycle \textit{Buffa-14} hypoxia metagene (base $\text{HR}^* = 1.132, p = 4.26 \times 10^{-6}$; +\textit{MKI67} $\text{HR}^* = 1.115, p = 7.80 \times 10^{-7}$). \textbf{(B)} Benjamini-Hochberg $-\log_{10}(q)$ versus post-\textit{MKI67} partial $R^2$ ($\%$ variance explained) across all $m = 286$ Sanger GDSC2 compounds evaluated on empirical $1 - \text{AUC}$, showing $0/286$ drugs below $q = 0.10$, highlighting the five top plateaued compounds (\texttt{WEHI-539}, \texttt{Navitoclax}, \texttt{Cisplatin}, \texttt{Elephantin}, and \texttt{ABT737}, all $q = 0.1126$), and noting the single small-sample outlier \texttt{SL0101} ($n = 110$, $\eta_p^2 = 4.76\%$, $q = 0.340$; baseline screen max $3.53\%$).}
\label{fig:summary}
\end{figure*}

Independent single-cell, pathway, proteomic, and genomic-instability datasets similarly showed weak or non-specific associations. In exploratory paired PDAC single-cell RNA-seq from \texttt{CRA001160} ($n = 23$ paired resections~\cite{peng2019,han2023}; BH-FDR across the 5 tested genes), pseudobulk expression confirmed higher malignant ductal vs.\ fibroblast expression of \textit{SLC7A1} ($+1.564\ \log_2\text{FC}$, raw $p = 2.38 \times 10^{-7}$, BH $q = 1.19 \times 10^{-6}$) and \textit{SMARCA4} ($+0.430$, raw $p = 0.005$, BH $q = 0.014$), whereas \textit{HIF1A} ($-0.217$, raw $p = 0.134$, BH $q = 0.168$) and gene-level \textit{SLC1A5} ($+0.275$, raw $p = 0.377$, BH $q = 0.377$) did not differ between tumor and fibroblast compartments and \textit{ABCB1} showed lower malignant expression at raw $p = 0.011$ ($-0.666$, BH $q = 0.019$, Bonferroni $p_{\text{adj}} = 0.055$). Across 50 MSigDB v2024.1 Hallmark gene sets~\cite{liberzon2015} tested with inter-gene correlation adjustment (competitive CAMERA~\cite{wu2012}), $0/50$ Hallmark sets cleared FDR $< 0.05$ in \texttt{TCGA-PAAD} ($\min q = 0.239$, despite 38/50 reaching unadjusted pre-ranked GSEA $q < 0.05$), whereas $10/50$ sets (dominated by G2M/mitotic and hormone-response sets) cleared CAMERA FDR in \texttt{METABRIC} (cross-cohort NES $\rho = +0.4182$, $p = 0.003$). Across 12 TCPA Level-4 RPPA cohorts~\cite{li2013} ($N = 4,192$), $98/194$ antibodies cleared FDR $q < 0.05$ owing to large sample size, with uniformly small effect sizes (max $|\rho| = 0.1583$ for ASNS; Bcl-2 $\rho = -0.1537$). Cohort-fixed-effect regression against Knijnenburg chromosomal instability scores~\cite{knijnenburg2018} showed negligible shared variance (partial $\Delta R^2 = 0.98\%$ across 7 cohorts, $N = 4,045$; $\Delta R^2 = 0.58\%$ across extended TCGA cohorts, $N = 5,604$ aneuploidy / $7,770$ HRD).

\section{Discussion}
Across clinical, functional genomic, perturbation, and pharmacogenomic datasets, an unweighted bulk \textit{SMARCA4} + \textit{HIF1A} mRNA score does not provide generalizable prognostic or therapeutic stratification in solid tumors. Across the 18-cohort TCGA panel ($N = 7,382$), the composite hazard ratio is statistically null before ($\text{HR}^* = 1.044$, $p = 0.162$) and after ($\text{HR}^* = 1.018$, $p = 0.545$) adjustment for \textit{MKI67}, matching the 19-cohort sensitivity pool ($\text{HR}^* = 1.020$, $p = 0.488$ post-\textit{MKI67}), with $0/18$ ($0/19$) cohorts clearing BH-FDR post-\textit{MKI67}. Across independent non-TCGA cohorts, \textit{MKI67} control abolished \texttt{CGGA-325} glioma stratification ($\text{HR} = 1.006$, $p = 0.943$), receptor/subtype adjustment attenuated \texttt{BRCA-METABRIC} ($p = 0.361$), and single-gene sensitivity checks across the 10-cohort clinical subset ($N = 5,102$) confirmed that \textit{HIF1A} mRNA alone was not associated with overall survival ($\text{HR}^* = 1.047$, $p = 0.222$).

The contrast with \textit{Buffa-14} shows why bulk \textit{HIF1A} mRNA fails as a hypoxia surrogate. Whereas \textit{HIF1A} mRNA does not rise during experimental $1\%\ \text{O}_2$ hypoxia across four independent GEO studies, the 14 non-cell-cycle downstream hypoxia genes in \textit{Buffa-14} capture hypoxia-response transcriptional activity, stratifying survival across the 18-cohort TCGA panel even after \textit{MKI67} control ($\text{HR}^* = 1.115$, $p = 7.80 \times 10^{-7}$, $I^2 = 0.9\%$) and improving omnibus nested model fit ($\chi^2(18) = 42.67$, $p = 8.90 \times 10^{-4}$). Clinical hypoxia stratification should therefore rely on validated downstream target signatures or protein-level HIF-1$\alpha$/HIF-2$\alpha$ stabilization assays rather than raw \textit{HIF1A} transcript abundance.

Correcting lineage mismatches and fitness-score interpretations in DepMap 24Q2 and Sanger GDSC2 leads to consistent conclusions. When Broad DepMap 24Q2 cell lines were restricted to the three lineages matching the TCGA expression panel (Breast, CNS/Brain, Pancreas), \textit{HIF1A}--\textit{SLC7A1} co-expression did not differ significantly from primary tumors ($Z = 1.02$, $p = 0.305$ across $n = 229$ unique models). In \textit{VHL}-LoF lines ($n = 19$), \textit{EPAS1} exhibits selective negative Chronos dependency ($-0.3464$) while \textit{HIF1A} exhibits positive/non-essential Chronos scores ($+0.3690$ vs.\ $+0.1127$ in WT), so the correlation between \textit{HIF1A} and \textit{SLC7A1} Chronos scores ($r = +0.5412$, $p = 0.017$) reflects co-variation among non-essential fitness scores rather than lethal hypoxic dependency. Likewise, applying \textit{MKI67} control across all $286$ GDSC2 compounds leaves $0/286$ FDR-significant drug hits ($q \ge 0.1126$).

Five limitations bound these conclusions. First, \textit{SMARCA4} is frequently inactivated by truncating mutations or focal deletions in solid tumors, so bulk mRNA abundance without somatic mutation or copy-number stratification is an incomplete proxy for functional SWI/SNF (BAF) complex activity. Second, proliferation control across the 18-cohort TCGA panel and GDSC2 screens relied primarily on single-gene \textit{MKI67} (with \textit{TOP2A} additionally checked in \texttt{CGGA-325}) rather than a multi-gene cell-cycle metagene. Third, multi-cohort TCGA Cox models did not adjust for patient sex or tumor purity outside \texttt{TCGA-PAAD}, \texttt{TCGA-KIRC}, and \texttt{TCGA-BLCA}. Fourth, GDSC2 $1 - \text{AUC}$ linear models lack cell-doubling time---a known confounder of cytotoxic agents such as Cisplatin---because doubling times are absent from DepMap 24Q2 model annotations. Fifth, across the 30 solid TCGA tumor types (19 in Table~\ref{tab:survival}, where \texttt{TCGA-COADREAD} merges \texttt{COAD} and \texttt{READ}, plus \texttt{TCGA-PRAD} in the 19-cohort sensitivity pool), 10 solid cohorts (\texttt{LGG}, \texttt{KIRP}, \texttt{KICH}, \texttt{MESO}, \texttt{UCS}, \texttt{UVM}, \texttt{ACC}, \texttt{CHOL}, \texttt{PCPG}, and \texttt{TGCT}) were excluded because they lacked complete cBioPortal PanCanAtlas/high-mortality RSEM and clinical archives in our TCGA pipeline; \texttt{TCGA-LGG} diffuse glioma biology was evaluated instead via the independent \texttt{CGGA-325} cohort ($N = 313$), while omitting \texttt{KIRP} and \texttt{KICH} leaves papillary and chromophobe renal carcinomas unassessed.

These negative findings point to four practical requirements for future studies of hypoxia and SWI/SNF signaling. First, clinical hypoxia stratification should rely on validated downstream transcriptional metagenes (such as non-cell-cycle \textit{Buffa-14}~\cite{buffa2010}) or direct protein-level HIF-1$\alpha$/HIF-2$\alpha$ stabilization assays rather than steady-state \textit{HIF1A} mRNA. Second, therapeutic targeting of the SWI/SNF complex should stratify tumors by somatic \textit{SMARCA4} truncating mutations or focal deletions---and exploit \textit{SMARCA2} paralog synthetic lethality~\cite{hoffman2014}---rather than bulk \textit{SMARCA4} transcript abundance. Third, evaluating metabolic adaptation via \textit{SLC1A5} requires isoform-resolved transcriptomics or mitochondrial sub-fractionation to separate hypoxia-inducible \texttt{SLC1A5\_var}~\cite{yoo2020} from canonical plasma-membrane \textit{SLC1A5}. Fourth, pharmacogenomic screens testing hypoxia-dependent drug vulnerabilities should be performed under controlled hypoxic tension ($1\%\ \text{O}_2$) rather than standard ambient $21\%\ \text{O}_2$ culture.

\section{Methods}

\subsection{Clinical Cohorts, Composite Score, and Survival Modeling}
The unweighted composite score $Z(z_{\textit{SMARCA4}} + z_{\textit{HIF1A}})$ (within-cohort mean $0$, SD $1$) was evaluated across survival and drug cohorts, while downstream metabolic and efflux genes (\textit{SLC7A1}, \textit{SLC1A5}, \textit{ABCB1}) were examined separately in DepMap, GEO, and PDAC scRNA-seq. Complete-case multivariable Cox models evaluated overall survival ($\text{OS} > 0$) adjusting for age and clinical stage or histological grade (\texttt{UCEC} and \texttt{CESC} used grade; \texttt{SARC} and \texttt{GBM} used age; \texttt{CGGA-325} additionally adjusted for \textit{IDH} mutation and \textit{MGMT} methylation; \texttt{TCGA-PAAD} additionally adjusted for ABSOLUTE purity, with a uniform \texttt{Age + Stage} sensitivity check at $N = 174$). The primary 18-cohort TCGA panel ($N = 7,382$, $2,352$ deaths; 16 stage- or grade-adjusted, \texttt{SARC} and \texttt{GBM} age-adjusted) covers the major adult solid cohorts with complete clinical and RSEM RNA-seq tables~\cite{tcga2018}, with un-staged \texttt{TCGA-PRAD} ($N = 493$, $10$ deaths) added in a 19-cohort sensitivity pool ($N = 7,875$). Across the 30 solid TCGA tumor types, 19 are represented in Table~\ref{tab:survival} (\texttt{TCGA-COADREAD} merges \texttt{COAD} and \texttt{READ}), 1 in sensitivity (\texttt{TCGA-PRAD}), and 10 (\texttt{LGG}, \texttt{KIRP}, \texttt{KICH}, \texttt{MESO}, \texttt{UCS}, \texttt{UVM}, \texttt{ACC}, \texttt{CHOL}, \texttt{PCPG}, and \texttt{TGCT}) were excluded due to absent PanCanAtlas/high-mortality RSEM and clinical archives (\texttt{LGG} diffuse glioma was instead evaluated via \texttt{CGGA-325}). Low-event cohorts (\texttt{TCGA-THCA}, 16 deaths; \texttt{TCGA-PRAD}, 10 deaths) are naturally downweighted by inverse-variance random-effects meta-analysis. Single-gene \textit{SMARCA4} and \textit{HIF1A} Cox decomposition covered the 10-cohort clinical subset ($N = 5,102$: \texttt{TCGA-KIRC}, \texttt{TCGA-BLCA}, \texttt{TCGA-GBM}, \texttt{TCGA-BRCA}, \texttt{TCGA-SARC}, \texttt{TCGA-SKCM}, \texttt{TCGA-LIHC}, and \texttt{TCGA-PAAD} [$N = 3,323$], plus \texttt{CGGA-325} and \texttt{BRCA-METABRIC}). \textit{MKI67} was not available in the processed \texttt{ICGC/PCAWG-PAAD} expression extract, so only \texttt{Age + Sex} and unadjusted log-rank models were evaluated there. Pooled hazard ratios used DerSimonian-Laird random-effects meta-analysis, non-cell-cycle \textit{Buffa-14} excluded \textit{CDKN3}~\cite{buffa2010}, and multiple testing across individual cohorts used Benjamini-Hochberg FDR within the primary 18-cohort family and separately within the 19-cohort sensitivity family. All analyses ran in Python 3.14.7 (\texttt{lifelines} 0.30.3 for Cox models, \texttt{statsmodels} 0.15.0 for Type II ANOVA and BH-FDR, \texttt{scipy} 1.18.1, \texttt{pandas} 2.3.3, \texttt{numpy} 2.5.4, and \texttt{matplotlib} 3.11.2, with competitive \texttt{CAMERA}~\cite{wu2012} implemented in Python following the Wu and Smyth inter-gene correlation variance-inflation formulation [not benchmarked directly against R \texttt{limma::camera}]) on $\log_2(\text{RSEM}+1)$ (TCGA), $\log_2(\text{TPM}+1)$ (DepMap 24Q2), normalized log-expression (GEO), and $\ln(\text{CP10K}+1)$ (\texttt{CRA001160}) scales. All statistical tests were two-sided.

\subsection{DepMap 24Q2 Functional Genomics and GEO Hypoxia Reanalysis}
Raw Broad DepMap Public 24Q2 CRISPR-Cas9 Chronos, RNA-seq, model metadata, and somatic mutation matrices~\cite{depmap2021} were filtered to exclude hematopoietic and non-cancerous models, yielding $1,038$ solid tumor CRISPR lines and $16$ non-malignant diploid controls ($14$ diploid and $974$ solid lines with matched RNA-seq for \textit{MKI67}-adjusted OLS; cell-doubling time is unavailable in DepMap model annotations). Genome-wide \textit{SMARCA4} co-essentiality was evaluated continuously across all $N = 1,038$ solid lines and $18,434$ genes using Pearson $r$ and Spearman $\rho$ with BH-FDR. We evaluated \textit{HIF1A}--\textit{SLC7A1} co-expression across \texttt{TCGA-BRCA}, \texttt{TCGA-PAAD}, and \texttt{TCGA-GBM} ($N = 1,384$)---spanning epithelial, desmoplastic, and necrotic microenvironments---and matched them to Broad DepMap 24Q2 default models in Breast, CNS/Brain, and Pancreas ($n = 229$ unique; $n = 221$ solid-filtered) using Fisher's $r$-to-$z$ test, reporting per-lineage correlations alongside the pooled estimate. Strict \textit{VHL}-LoF models ($n = 19$ vs.\ $1,019$ WT) were filtered for truncating, splice, or predicted damaging variants. GEO $1\%$ vs.\ $21\%\ \text{O}_2$ comparisons (\texttt{GSE139673}~\cite{gse139673}, \texttt{GSE145108}~\cite{gse145108}, \texttt{GSE283995}~\cite{gse283995}, \texttt{GSE300765}~\cite{gse300765}) retained formal $t$-test $p$-values only for the 24 powered single-lineage comparisons ($n \ge 3$) with exploratory Benjamini-Hochberg FDR control across those 24 tests, reporting \texttt{GSE145108} ($n = 2$ pairs per lineage across RKO, A549, and H460) strictly descriptively without cross-organ pooling.

\subsection{Sanger GDSC2 Drug Screening and Multi-Omics Analyses}
Lineage-controlled Type II ANOVA models across $m = 286$ Sanger GDSC2 compounds~\cite{iorio2016} ($n = 110\text{--}597$ solid cell lines per drug, median $n = 547$; \texttt{SL0101} tested at $n = 110$) using bounded empirical $1 - \text{AUC}$ were adjusted with Benjamini-Hochberg step-up FDR across all 286 drugs after including standardized DepMap \textit{MKI67} RNA-seq. In exploratory PDAC scRNA-seq (\texttt{CRA001160}~\cite{peng2019,han2023}; $57,443$ total cells, including $41,976$ tumor-sample cells across $24$ patients and $n = 23$ paired resections containing both malignant ductal and fibroblast compartments), $\ln(\text{CP10K}+1)$ expression (which aggregates canonical \textit{SLC1A5} and \texttt{SLC1A5\_var} at the gene level) was back-transformed via $\exp(y) - 1$, averaged per patient within each compartment, compared via paired pseudobulk $\log_2\text{FC} = \log_2((\overline{\text{Mal}}_p + 0.01) / (\overline{\text{CAF}}_p + 0.01))$, and adjusted across the 5 tested genes using Benjamini-Hochberg FDR (with Bonferroni $p_{\text{adj}}$ noted for \textit{ABCB1}). MSigDB v2024.1 Hallmark CAMERA/GSEA~\cite{liberzon2015,wu2012}, 12-cohort TCPA Level-4 RPPA Spearman correlations~\cite{li2013}, and Knijnenburg chromosomal instability partial $\Delta R^2$ models~\cite{knijnenburg2018} were evaluated on the same unweighted two-gene score.

\subsection*{Data and Code Availability}
Analysis scripts, processed summary tables, and manuscript source files are available at \url{https://github.com/SachsinSinnathurai/smarca4-hif1a-pharmacogenomics}. Public datasets analyzed in this study are available from the TCGA PanCanAtlas Clinical Data Resource~\cite{tcga2018}, Broad DepMap Public 24Q2 (\url{https://depmap.org}), Sanger GDSC2 (\url{https://www.cancerrxgene.org}), NCBI Gene Expression Omnibus (\texttt{GSE139673}, \texttt{GSE145108}, \texttt{GSE283995}, and \texttt{GSE300765}), Genome Sequence Archive (\texttt{CRA001160} via TISCH2~\cite{han2023}), MSigDB v2024.1~\cite{liberzon2015}, and TCPA Level-4 RPPA~\cite{li2013}. Controlled-access clinical datasets (\texttt{BRCA-METABRIC}~\cite{curtis2012} and \texttt{ICGC/PCAWG-PAAD}~\cite{pcawg2020}) were analyzed under their respective consortium data-access terms.

\subsection*{Competing Interests}
The author declares no competing interests.

\begin{thebibliography}{26}
\bibitem{ivan2001} Ivan M, Kondo K, Yang H, et al. HIF$\alpha$ targeted for VHL-mediated destruction by proline hydroxylation: implications for O$_2$ sensing. \textit{Science}. 2001;292(5516):464--468. doi:10.1126/science.1059817 (PMID: 11292862)
\bibitem{jaakkola2001} Jaakkola P, Mole DR, Tian YM, et al. Targeting of HIF-$\alpha$ to the von Hippel-Lindau ubiquitylation complex by O$_2$-regulated prolyl hydroxylation. \textit{Science}. 2001;292(5516):468--472. doi:10.1126/science.1059796 (PMID: 11292861)
\bibitem{kenneth2009} Kenneth NS, Mudie S, van Uden P, Rocha S. SWI/SNF regulates the cellular response to hypoxia. \textit{J Biol Chem}. 2009;284(7):4123--4131. doi:10.1074/jbc.M808491200 (PMID: 19091745)
\bibitem{sena2013} Sena JA, Wang L, Hu CJ. BRG1 and BRM chromatin-remodeling complexes regulate the hypoxia response by acting as coactivators for a subset of hypoxia-inducible transcription factor target genes. \textit{Mol Cell Biol}. 2013;33(19):3849--3863. doi:10.1128/MCB.00731-13 (PMID: 23897425)
\bibitem{elorza2012} Elorza A, Soro-Arnaiz I, Melendez-Rodriguez F, et al. HIF2$\alpha$ acts as an mTORC1 activator through the amino acid carrier SLC7A5. \textit{Mol Cell}. 2012;48(5):681--691. doi:10.1016/j.molcel.2012.09.017 (PMID: 23063528)
\bibitem{yoo2020} Yoo HC, Park SJ, Nam M, et al. A variant of SLC1A5 is a mitochondrial glutamine transporter for metabolic reprogramming in cancer cells. \textit{Cell Metab}. 2020;31(2):267--283. doi:10.1016/j.cmet.2019.11.020 (PMID: 31866442)
\bibitem{buffa2010} Buffa FM, Harris AL, West CM, Miller CJ. Large meta-analysis of multiple cancers reveals a common, compact and highly prognostic hypoxia metagene. \textit{Br J Cancer}. 2010;102(2):428--435. doi:10.1038/sj.bjc.6605450 (PMID: 20087356)
\bibitem{gse139673} Markolin P, Davidson N, Hirt CK, et al. Identification of HIF-dependent alternative splicing in gastrointestinal cancers and characterization of a long, coding isoform of SLC35A3. \textit{Genomics}. 2021;113(2):515--529. doi:10.1016/j.ygeno.2020.12.039 (PMID: 33418078; NCBI GEO Series GSE139673: AsPC-1, siControl, 1\% vs.\ 21\% $\text{O}_2$, 24~h; $n = 3$ pairs).
\bibitem{gse145108} Andrysik Z, Bender H, Galbraith MD, Espinosa JM. Multi-omics analysis reveals contextual tumor suppressive and oncogenic gene modules within the acute hypoxic response. \textit{Nat Commun}. 2021;12(1):1375. doi:10.1038/s41467-021-21687-2 (PMID: 33654095; NCBI GEO Series GSE145108: RKO, A549, and H460 under 1\% vs.\ 21\% $\text{O}_2$, 24~h; $n = 2$ pairs per cell line, descriptive).
\bibitem{gse283995} Li L, Hammerlindl H, Shen SQ, Bao F, Hammerlindl S, Altschuler SJ, Wu LF. A phenopushing platform to identify compounds that alleviate acute hypoxic stress by fast-tracking cellular adaptation. \textit{Nat Commun}. 2025;16(1):2684. doi:10.1038/s41467-025-57754-1 (PMID: 40102413; NCBI GEO Series GSE283995: HepG2 cells under 1\% vs.\ 21\% $\text{O}_2$ at 24~h [$n = 6$ vs.\ $6$] and day 6 [$n = 5$ vs.\ $6$]).
\bibitem{gse300765} B\aa ng-Rudenstam A, Cerezo-Maga\~na M, Horv\'ath M, et al. Tumour acidosis remodels the glycocalyx to control lipid scavenging and ferroptosis. \textit{Nat Cell Biol}. 2026;28(3):567--580. doi:10.1038/s41556-026-01879-y (PMID: 41673170; NCBI GEO Series GSE300765: U87MG cells under 1\% vs.\ 21\% $\text{O}_2$, 48~h; $n = 3$ pairs).
\bibitem{tcga2018} Liu J, Lichtenberg T, Hoadley KA, et al. An Integrated TCGA Pan-Cancer Clinical Data Resource to Drive High-Quality Survival Outcome Analytics. \textit{Cell}. 2018;173(2):400--416. doi:10.1016/j.cell.2018.02.052 (PMID: 29625055)
\bibitem{cgga2021} Zhao Z, Zhang KN, Wang Q, et al. Chinese Glioma Genome Atlas (CGGA): A Comprehensive Resource with Functional Genomic Data from Chinese Glioma Patients. \textit{Genomics Proteomics Bioinformatics}. 2021;19(1):1--12. doi:10.1016/j.gpb.2020.10.005 (PMID: 33662628)
\bibitem{curtis2012} Curtis C, Shah SP, Chin SF, et al. The genomic and transcriptomic architecture of 2,000 breast tumours reveals novel subgroups. \textit{Nature}. 2012;486(7403):346--352. doi:10.1038/nature10983 (PMID: 22522925)
\bibitem{pcawg2020} ICGC/TCGA Pan-Cancer Analysis of Whole Genomes Consortium. Pan-cancer analysis of whole genomes. \textit{Nature}. 2020;578(7793):82--93. doi:10.1038/s41586-020-1969-6 (PMID: 32025007)
\bibitem{carter2012} Carter SL, Cibulskis K, Helman E, et al. Absolute quantification of somatic DNA alterations in human cancer. \textit{Nat Biotechnol}. 2012;30(5):413--421. doi:10.1038/nbt.2203 (PMID: 22544022)
\bibitem{raphael2017} Raphael BJ, Hruban RH, Aguirre AJ, et al. Integrated Genomic Characterization of Pancreatic Ductal Adenocarcinoma. \textit{Cancer Cell}. 2017;32(2):185--203. doi:10.1016/j.ccell.2017.07.007 (PMID: 28810144)
\bibitem{depmap2021} Dempster JM, Boyle I, Vazquez F, et al. Chronos: a cell population dynamics model of CRISPR experiments that improves inference of gene fitness effects. \textit{Genome Biol}. 2021;22(1):343. doi:10.1186/s13059-021-02540-7 (PMID: 34930405)
\bibitem{iorio2016} Iorio F, Knijnenburg TA, Vis DJ, et al. A Landscape of Pharmacogenomic Interactions in Cancer. \textit{Cell}. 2016;166(3):740--754. doi:10.1016/j.cell.2016.06.017 (PMID: 27397505); together with the Sanger GDSC2 public release (\url{https://www.cancerrxgene.org}).
\bibitem{peng2019} Peng J, Sun BF, Chen CY, et al. Single-cell RNA-seq highlights intra-tumoral heterogeneity and malignant progression in pancreatic ductal adenocarcinoma. \textit{Cell Res}. 2019;29(9):725--738. doi:10.1038/s41422-019-0195-y (PMID: 31273297)
\bibitem{han2023} Han Y, Wang Y, Dong X, et al. TISCH2: expanded datasets and new tools for single-cell transcriptome analyses of the tumor microenvironment. \textit{Nucleic Acids Res}. 2023;51(D1):D1425--D1431. doi:10.1093/nar/gkac959 (PMID: 36321662)
\bibitem{liberzon2015} Liberzon A, Birger C, Thorvaldsdottir H, Ghandi M, Mesirov JP, Tamayo P. The Molecular Signatures Database (MSigDB) hallmark gene set collection. \textit{Cell Syst}. 2015;1(6):417--425. doi:10.1016/j.cels.2015.12.004 (PMID: 26771021)
\bibitem{wu2012} Wu D, Smyth GK. Camera: a competitive gene set test accounting for inter-gene correlation. \textit{Nucleic Acids Res}. 2012;40(17):e133. doi:10.1093/nar/gks461 (PMID: 22638577)
\bibitem{li2013} Li J, Lu Y, Akbani R, et al. TCPA: a resource for cancer functional proteomics data. \textit{Nat Methods}. 2013;10(11):1046--1047. doi:10.1038/nmeth.2650 (PMID: 24037243)
\bibitem{knijnenburg2018} Knijnenburg TA, Wang L, Zimmermann MT, et al. Genomic and molecular landscape of DNA damage repair deficiency across The Cancer Genome Atlas. \textit{Cell Rep}. 2018;23(1):239--254. doi:10.1016/j.celrep.2018.03.076 (PMID: 29617664)
\bibitem{hoffman2014} Hoffman GR, Rahal R, Buxton F, et al. Functional epigenetics approach identifies BRM/SMARCA2 as a critical synthetic lethal target in BRG1-deficient cancers. \textit{Proc Natl Acad Sci U S A}. 2014;111(8):3128--3133. doi:10.1073/pnas.1316793111 (PMID: 24520176)
\end{thebibliography}

\end{document}
